\section{Contacts between actual fine cells}
\label{sec:area_law_cell_contacts}

Fix an origin $o$, finite endpoint set $Z$ and nonnegative integer width
$C_0$. Retain the fine exponents $\ell_k$ and index sets $F_{k,\ell_k}$,
and write $Q_{k,z}=C_{\ell_k,z}$ for an actual fine cell. The common
dyadic origin restricts positive-length contacts to whole sides and
midpoint halves; see \cite[Section~11]{OpenAI2026AreaLaw}.

% Source: 10-geometry.tex, lines 173--181, 299--306 and 352--356.
% The statements below concern indexed fine cells and optional midpoint
% halves. Matching under the actual mask and opposing-region labels
% remain further obligations; the full two-family proposition is not tagged.

\begin{definition}[Whole sides of a dyadic cell]
    \label{def:area_law_dyadic_cell_side}
    \lean{TNLean.PEPS.AreaLaw.Geometry.dyadicCellSide}
    \leanok
    \uses{def:area_law_cell_fan}
    For every exponent $\ell\ge0$ and cell index $z$, number the
    coordinate sides of $\overline{C_{\ell,z}}$ counterclockwise,
    starting with the right side. If $a_{\ell,z,s}$ and $b_{\ell,z,s}$
    are the corresponding unsplit fan endpoints, put
    $S_{\ell,z,s}=[a_{\ell,z,s},b_{\ell,z,s}]$.
    Side indices are read modulo four. Writing $a=a_{\ell,z,s}$ and
    $b=b_{\ell,z,s}$, its midpoint halves are
    $H_{\ell,z,s,0}=[a,(a+b)/2]$ and
    $H_{\ell,z,s,1}=[(a+b)/2,b]$.
\end{definition}

\begin{theorem}[Distinct actual fine cells are disjoint]
    \label{thm:area_law_fine_cells_disjoint}
    \lean{TNLean.PEPS.AreaLaw.Geometry.fineLayer_cells_disjoint}
    \leanok
    \uses{def:area_law_fine_layer_indices,def:area_law_fine_belt_scales}
    For arbitrary $C_0\ge0$ and layer indices $k,h$, suppose that
    $z\in F_{k,\ell_k}$, $w\in F_{h,\ell_h}$ and $(k,z)\ne(h,w)$.
    Then $Q_{k,z}\cap Q_{h,w}=\varnothing$.
\end{theorem}
\begin{proof}\leanok
    \uses{thm:area_law_dyadic_indices,thm:area_law_layer_pairwise_disjoint,
        thm:area_law_fine_cell_containment,thm:area_law_fine_belt_scales}
    If $k=h$, the unique cell-index criterion implies that any common
    point would give $z=w$. If $k\ne h$, the two fine cells belong to
    the disjoint half-open layers $D_k$ and $D_h$, since
    $\ell_k\le k$ and $\ell_h\le h$.
\end{proof}

\begin{theorem}[Positive-length contact between fine cells]
    \label{thm:area_law_fine_cell_contact}
    \lean{TNLean.PEPS.AreaLaw.Geometry.fineLayer_closedCell_contact}
    \leanok
    \uses{def:area_law_dyadic_cell_side,def:area_law_fine_layer_indices,
        def:area_law_fine_belt_scales}
    Suppose that $C_0\ge2$, $k\ge50\,000\,000$, $z\in F_{k,\ell_k}$,
    $w\in F_{h,\ell_h}$ and $(k,z)\ne(h,w)$. If
    $I=\overline{Q_{k,z}}\cap\overline{Q_{h,w}}$ contains two distinct
    points, there are opposite side indices $s,t$, with $t=s+2$ modulo
    four, for which one of the following alternatives holds:
    \begin{enumerate}
        \item $\ell_k=\ell_h$ and
              $I=S_{\ell_k,z,s}=S_{\ell_h,w,t}$.
        \item $\ell_h=\ell_k+1$ and
              $I=S_{\ell_k,z,s}=H_{\ell_h,w,t,n}$ for some $n\in\{0,1\}$.
        \item $\ell_k=\ell_h+1$ and
              $I=S_{\ell_h,w,t}=H_{\ell_k,z,s,n}$ for some $n\in\{0,1\}$.
    \end{enumerate}
    The opposing index $h$ is unrestricted.
\end{theorem}
\begin{proof}\leanok
    \uses{thm:area_law_fine_cells_disjoint,thm:area_law_fine_cell_containment,
        thm:area_law_local_layer_indices,thm:area_law_adjacent_fine_meshes,
        thm:area_law_dyadic_closures,def:area_law_dyadic_cell_side}
    A common point belongs to both closed layers. The nearby-layer
    estimate gives $h\le k+1$ and $k\le h+1$, so the fine exponents
    differ by at most one. Order the cells by their fine exponents.
    In units of the smaller side, translate the larger cell's lower
    corner to the origin. The two half-open cells are then
    $[a,a+1)\times[b,b+1)$ and $[0,m)^2$, where $a,b$ are integers
    and $m\in\{1,2\}$.

    Closed contact gives $-1\le a,b\le m$. Disjointness excludes
    $0\le a,b<m$, and two distinct common points exclude the case
    where both coordinates lie in $\{-1,m\}$. Thus one coordinate
    equals $-1$ or $m$, while the other is an integer in $[0,m)$.
    The common set is therefore a whole smaller side on the opposite
    coordinate side of the larger square. For $m=1$ it is the whole
    larger side; for $m=2$ it is one of its midpoint halves. Swapping
    the two cells gives the third alternative.
\end{proof}
